\documentclass[a4paper,12pt]{article}

\usepackage[utf8]{inputenc}
\usepackage{geometry}
\usepackage{graphicx}
\graphicspath{{Figures/}}
\usepackage{mathptmx}
\usepackage{setspace}
\usepackage{titlesec}
\usepackage{parskip}
\usepackage{enumitem}
\usepackage{indentfirst}
\usepackage{listings}
\usepackage{xcolor}
\usepackage{hyperref}
\usepackage{xurl}
\usepackage{microtype}
\usepackage{float}
\setlength{\emergencystretch}{1em}
\sloppy

\lstset{
    breaklines=true,
    postbreak=\mbox{\textcolor{red}{$\hookrightarrow$}\space},
}

\setlength{\parindent}{1em}

\definecolor{dkgreen}{rgb}{0,0.6,0}
\definecolor{gray}{rgb}{0.5,0.5,0.5}
\definecolor{mauve}{rgb}{0.58,0,0.82}
\definecolor{dkblue}{rgb}{0,0,0.5}
\definecolor{lightblue}{rgb}{0,0.5,0.8}
\definecolor{darkred}{rgb}{0.6,0,0}
\definecolor{backcolour}{rgb}{0.97,0.97,0.97}

\lstdefinestyle{javastyle}{
    language=Java,
    backgroundcolor=\color{backcolour},   
    commentstyle=\color{dkgreen}\itshape,
    keywordstyle=\color{dkblue}\bfseries,
    stringstyle=\color{mauve},
    basicstyle=\ttfamily\footnotesize,
    breakatwhitespace=false,         
    breaklines=true,                 
    captionpos=b,                    
    keepspaces=true,                 
    numbers=left,                    
    numbersep=5pt,                  
    showspaces=false,                
    showstringspaces=false,
    showtabs=false,                  
    tabsize=4,
    frame=lines,
    rulecolor=\color{gray},
    moredelim=[s][\color{orange}]{@}{\ },
}

\lstdefinestyle{xmlstyle}{
    language=XML,
    backgroundcolor=\color{backcolour},
    basicstyle=\ttfamily\footnotesize,
    morekeywords={xmlns,version,type,scope,groupId,artifactId,dependencies,dependency,project,modelVersion,build,plugins,plugin,configuration,source,target},
    keywordstyle=\color{darkred}\bfseries,
    commentstyle=\color{dkgreen}\itshape,
    stringstyle=\color{dkblue},
    tagstyle=\color{darkred},
    breaklines=true,
    numbers=left,
    numbersep=5pt,
    tabsize=2,
    frame=lines
}

\geometry{
    top=3cm,
    left=4cm,
    right=3cm,
    bottom=3cm
}

\renewcommand{\thesection}{\Alph{section}} 

\renewcommand{\thesubsection}{\thesection.\arabic{subsection}}

\titleformat{\section}
  {\normalfont\Large\bfseries}
  {\thesection.}
  {1em}
  {}

\begin{document}

\begin{titlepage}
    \centering
    \onehalfspacing

    {\fontsize{14}{17}\selectfont \textbf{LAPORAN PRAKTIKUM}} \\[0.2cm]
    {\fontsize{16}{19}\selectfont \textbf{PRAKTIKUM PENGUJIAN PERANGKAT LUNAK}} \\[0.5cm]
    {\fontsize{14}{17}\selectfont \textbf{PERTEMUAN 3}} \\[0.2cm]
    {\fontsize{14}{17}\selectfont \textbf{Parameterized Testing}}
    
    \vfill
    
\includegraphics[width=6cm]{Figures/UGM.png} 
    \vfill
    
    {\large Disusun Oleh:} \\[0.5cm]
    
    \begin{tabular}{l c l} 
        Nama & : & Abdullah Afif Habiburrohman \\
        NIM  & : & 537611/SV/24441 \\
        Kelas & : & BB / B2 \\
        Dosen & : & Divi Galih Prasetyo Putri, S.Kom., M.Kom., Ph.D.
    \end{tabular}
    
    \vfill
    
    {\fontsize{12}{14}\selectfont \textbf{
        PROGRAM STUDI SARJANA TERAPAN \\ 
        TEKNOLOGI REKAYASA PERANGKAT LUNAK \\
        DEPARTEMEN TEKNIK ELEKTRO DAN INFORMATIKA \\
        SEKOLAH VOKASI \\
        UNIVERSITAS GADJAH MADA \\
        YOGYAKARTA \\
        2026
    }}
    
\end{titlepage}

\tableofcontents
\newpage

\section{Tujuan Pembelajaran}
\begin{enumerate}
    \item Mahasiswa mampu memahami tipe-tipe parameter dalam test method
    \item Mahasiswa mampu megnimplementasikan tipe-tipe parameter pada pengujian
\end{enumerate}

\section{Dasar Teori}
Parameterized test adalah fitur dalam ekosistem JUnit 5 yang memungkinkan eksekusi sebuah metode pengujian secara berulang kali menggunakan serangkaian nilai input atau parameter yang berbeda. Fitur ini memfasilitasi pengujian untuk berbagai kasus dan skenario tanpa mengharuskan pengembang menulis metode pengujian terpisah untuk masing-masing input.

Penerapan pengujian terparameterisasi sangat erat kaitannya dengan prinsip perangkat lunak DRY (Don't Repeat Yourself), di mana redundansi kode dapat ditekan dengan menghindari penyalinan kode pengujian yang identik secara berulang. Selain menjaga agar basis kode pengujian tetap bersih dan mudah dibaca, pendekatan ini juga secara efektif meningkatkan cakupan pengujian (test coverage) untuk berbagai skenario khusus tanpa mengorbankan kemudahan pemeliharaan.

Untuk mengimplementasikan fitur ini, pengujian tidak lagi ditandai dengan anotasi standar, melainkan harus menggunakan anotasi @ParameterizedTest. Secara arsitektural, implementasi pengujian ini membutuhkan penambahan dependensi junit-jupiter-params pada konfigurasi proyek, dan setiap metode pengujian terparameterisasi diwajibkan untuk mendeklarasikan minimal satu sumber argumen (source of arguments).

JUnit Jupiter menyediakan berbagai macam anotasi sumber untuk menyuplai nilai parameter ke dalam metode pengujian. Beberapa mekanisme utama penyediaan parameter tersebut meliputi:

\begin{enumerate}
    \item Anotasi @ValueSource\\
    Anotasi ini digunakan untuk mendefinisikan array nilai literal tunggal (single-value parameter). Tipe data yang didukung mencakup tipe primitif (seperti int, long, double), String, maupun Class. JUnit 5 akan menjalankan metode pengujian sebanyak jumlah elemen yang ada di dalam array tersebut, dengan menyuntikkan satu elemen pada setiap eksekusi pengujian. Namun, tipe sumber argumen ini tidak mendukung penyisipan argumen yang bernilai null.

    \item Anotasi @CsvSource dan @CsvFileSource\\
    Untuk skenario di mana metode pengujian memerlukan penyisipan lebih dari satu argumen sekaligus (misalnya sepasang nilai input awal dan ekspektasi nilai akhir), sumber data berformat CSV (Comma-Separated Values) dapat dimanfaatkan.
    \begin{itemize}
        \item @CsvSource memungkinkan deklarasi deretan parameter secara langsung dalam bentuk sekumpulan string yang dipisahkan oleh tanda koma di dalam anotasi kode sumber.
        \item @CsvFileSource beroperasi serupa, namun argumen dimuat secara eksternal melalui file CSV yang tersimpan di dalam struktur repositori proyek, sehingga memungkinkan pemisahan yang rapi antara data pengujian dan logika kode.
    \end{itemize}

    \item Anotasi @MethodSource\\
    Apabila pengujian membutuhkan pasokan argumen yang direpresentasikan oleh objek yang lebih kompleks, parameter dapat disuplai melalui pemanggilan sebuah metode penyedia (factory method). Metode penyedia parameter ini diwajibkan untuk bersifat statis (static method) dan harus mengembalikan kumpulan argumen dalam format struktur data iteratif, seperti Stream, Iterable, Collection, ataupun sekumpulan array objek.
\end{enumerate}

\section{Hasil dan Pembahasan}
Mahasiswa melanjutkan praktikum sebelumnya, Wallet dan Testing Lifecycle pada kelas WalletTest. Mahasiswa akan menampilkan kode tambahan dan perubahan kode untuk mengimplementasikan tugas-tugas dalam parameterized test. Berikut adalah kode-kode yang telah diimplementasikan untuk memenuhi tugas praktikum ini:

\subsection{Kode Baru}
Berikut adalah kode-kode baru yang sebelumnya tidak ada pada praktikum sebelumnya:

\begin{lstlisting}[style=javastyle, caption=InsufficientFundsException.java]
package org.example;

public class InsufficientFundsException extends Exception {
    public InsufficientFundsException(String message) {
        super(message);
    }
}
\end{lstlisting}

Kode diatas merupakan penambahan exception baru yang akan digunakan untuk mengindikasikan bahwa saldo tidak mencukupi untuk melakukan transaksi tertentu.

\begin{lstlisting}[style=javastyle, caption=Owner.java]
package org.example;

public class Owner {
    private final String id;
    private final String name;
    private final String email;

    public Owner(String id, String name, String email) {
        this.id = id;
        this.name = name;
        this.email = email;
    }

    public String getId() { return id; }
    public String getName() { return name; }
    public String getEmail() { return email; }
}
\end{lstlisting}

Kode diatas adalah kode baru untuk kelas Owner yang sebelumnya hanya string pada kelas Wallet. Kelas Owner ini akan digunakan untuk menyimpan informasi pemilik dompet, seperti id, nama, dan email.

\subsection{Perubahan dan penambahan kode}
Berikut adalah kode-kode yang telah diubah dan ditambahkan untuk mengimplementasikan tugas-tugas dalam parameterized test:
\begin{lstlisting}[style=javastyle, caption=Card.java]
package org.example;

public class Card {
    private String namaBank;
    private int nomorRekening;

    public Card(String namaBank, int nomorRekening) {
        this.namaBank = namaBank;
        this.nomorRekening = nomorRekening;
    }

    public String getNamaBank() {
        return namaBank;
    }

    public int getNomorRekening() {
        return nomorRekening;
    }
}
\end{lstlisting}

Kode yang diubah adalah tipe data dari variabel nomorRekening yang sebelumnya berupa String menjadi int. Perubahan ini dilakukan untuk menyesuaikan dengan format nomor rekening yang umumnya berupa angka.

\begin{lstlisting}[style=javastyle, caption=Wallet.java]
package org.example;

import java.util.ArrayList;
import java.util.List;

public class Wallet {
    private Owner owner;
    private List<Card> cards;
    private double cash;

    public Wallet(Owner owner) {
        this.owner = owner;
        this.cards = new ArrayList<>();
        this.cash = 0.0;
    }

    public void setOwner(Owner owner) {
        this.owner = owner;
    }

    public Owner getOwner() {
        return owner;
    }

    public void addCards(String bank, int nomor) {
        Card newCard = new Card(bank, nomor);
        this.cards.add(newCard);
    }

    public Card removeCard(int nomor) {
        for (Card card : cards) {
            if (card.getNomorRekening() == nomor) {
                cards.remove(card);
                return card;
            }
        }
        return null;
    }

    public List<Card> getCards() {
        return cards;
    }

    public void deposit(double amount) {
        if (amount > 0) {
            this.cash += amount;
        }
    }

    public void withdraw(double amount) throws InsufficientFundsException {
        if (amount <= 0) {
            throw new IllegalArgumentException("Withdraw amount must be > 0");
        }
        if (this.cash < amount) {
            throw new InsufficientFundsException("Insufficient funds");
        }
        this.cash -= amount;
    }

    public double getCash() {
        return cash;
    }

    public void cleanCash() {
        this.cash = 0.0;
    }

    public void cleanCards() {
        cards.clear();
    }
}
\end{lstlisting}
Kode diatas adalah kode yang telah diubah dan ditambahkan pada kelas Wallet. Perubahan utama adalah:
\begin{enumerate}
    \item Mengubah owner dan turunannya dari tipe String menjadi tipe Owner untuk menyimpan informasi pemilik dompet secara lebih terstruktur.
    \item Mengubah tipe data variabel cash menjadi double untuk memungkinkan penyimpanan nilai uang dengan desimal.
    \item Mengimplementasikan perubahan tipe data variabel nomorRekening dari String menjadi int pada kelas Card ke kelas Wallet.
    \item Mengubah metode withdraw menjadi melemparkan InsufficientFundsException jika jumlah penarikan melebihi saldo yang tersedia, serta menambahkan validasi untuk memastikan jumlah penarikan harus lebih besar dari 0.
\end{enumerate}

\begin{lstlisting}[style=javastyle, caption=WalletTest.java]
package org.example;

import org.junit.jupiter.api.*;
import org.junit.jupiter.params.ParameterizedTest;
import org.junit.jupiter.params.provider.Arguments;
import org.junit.jupiter.params.provider.CsvFileSource;
import org.junit.jupiter.params.provider.MethodSource;
import org.junit.jupiter.params.provider.ValueSource;

import java.util.stream.Stream;

import static org.junit.jupiter.api.Assertions.*;

@TestInstance(TestInstance.Lifecycle.PER_CLASS)
@TestMethodOrder(MethodOrderer.OrderAnnotation.class)
class WalletTest {

    Wallet dompet;

    @BeforeAll
    void initClass() {
        Owner defaultOwner = new Owner("1", "Abdullah", "abd@ugm.ac.id");
        dompet = new Wallet(defaultOwner);
    }

    @AfterAll
    void cleanClass() {
        dompet = null;
    }

    @BeforeEach
    void initMethod() {
        dompet.deposit(10000.0);
        dompet.addCards("DefaultBank", 0);
    }

    @AfterEach
    void cleanMethod() {
        dompet.cleanCards();
        dompet.cleanCash();
    }

    @Test
    @Order(1)
    @DisplayName("Execute constructor, verify fields")
    void testIsWallet() {
        assertNotNull(dompet, "Objek dompet tidak boleh null");
        assertEquals("Abdullah", dompet.getOwner().getName());
    }

    @ParameterizedTest
    @ValueSource(doubles = {10000.0, 50000.0, 100000.0})
    @DisplayName("Tugas 1: Test Nominal Cash Valid (Deposit)")
    void testCashValid(double nominal) {
        dompet.cleanCash();
        dompet.deposit(nominal);
        assertEquals(nominal, dompet.getCash());
    }

    @ParameterizedTest
    @ValueSource(doubles = {-1000.0, -5000.0, 0.0})
    @DisplayName("Tugas 1: Test Nominal Cash Tidak Valid (Deposit diabaikan)")
    void testCashInvalid(double nominal) {
        dompet.cleanCash();
        dompet.deposit(nominal);
        assertEquals(0.0, dompet.getCash()); // Saldo tetap 0
    }

    @ParameterizedTest
    @CsvFileSource(resources = "/valid-withdraw.csv", numLinesToSkip = 1)
    @DisplayName("Tugas 2: Test Withdraw Valid (CSV)")
    void testValidWithdrawCSV(double deposit, double withdraw, double expectedTotal) throws InsufficientFundsException {
        dompet.cleanCash();
        dompet.deposit(deposit);

        if (withdraw > 0) {
            dompet.withdraw(withdraw);
        }

        assertEquals(expectedTotal, dompet.getCash());
    }

    @ParameterizedTest
    @CsvFileSource(resources = "/invalid-withdraw.csv", numLinesToSkip = 1)
    @DisplayName("Tugas 2: Test Withdraw Invalid Exception (CSV)")
    void testInvalidWithdrawCSV(double deposit, double withdraw, String exceptionType) {
        dompet.cleanCash();
        dompet.deposit(deposit);

        if (exceptionType.equals("InsufficientFundsException")) {
            assertThrows(InsufficientFundsException.class, () -> dompet.withdraw(withdraw));
        } else if (exceptionType.equals("IllegalArgumentException")) {
            assertThrows(IllegalArgumentException.class, () -> dompet.withdraw(withdraw));
        }
    }

    static Stream<Arguments> provideOwnerObjects() {
        return Stream.of(
                Arguments.of(new Owner("1", "Agus", "agus@ugm.ac.id")),
                Arguments.of(new Owner("2", "Siti", "siti@ugm.ac.id"))
        );
    }

    @ParameterizedTest
    @MethodSource("provideOwnerObjects")
    @DisplayName("Tugas 3: Test Set Owner Menggunakan Object")
    void testSetOwnerObject(Owner ownerParam) {
        dompet.setOwner(ownerParam);
        assertNotNull(dompet.getOwner());
        assertEquals(ownerParam.getName(), dompet.getOwner().getName());
        assertEquals(ownerParam.getEmail(), dompet.getOwner().getEmail());
    }
}
\end{lstlisting}
Kode diatas adalah kode yang telah diubah dan ditambahkan pada kelas WalletTest. Perubahan utama adalah:
\begin{enumerate}
    \item Mengimplementasikan perubahan pada class Wallet yang telah dilakukan sebelumnya.
    \item Menambahkan @ParameterizedTest pada metode testCashValid dan testCashInvalid dengan value source berupa array double untuk menguji berbagai nominal deposit yang valid dan tidak valid.
    \item Menambahkan @ParameterizedTest pada metode testValidWithdrawCSV dan testInvalidWithdrawCSV dengan value source berupa file CSV untuk menguji berbagai kasus withdraw yang valid dan tidak valid.
    \item Menambahkan @ParameterizedTest pada metode testSetOwnerObject dengan value source berupa objek Owner untuk menguji berbagai kasus owner yang valid.
\end{enumerate}

\subsection{Gambar Hasil Run}
Berikut adalah gambar hasil run dari kode diatas:

\begin{figure}[H]
    \centering
    \includegraphics[width=1\textwidth]{one.png}
    \caption{Hasil Run}
    \label{fig:hasil_run}
\end{figure}

\subsection{GitHub}
Berikut adalah link repository di GitHub yang berisi kode diatas \url{https://github.com/bukunya/PPPL/tree/Pertemuan3}

\section{Kesimpulan}
Pada praktikum ini mahasiswa telah mempelajari dan mengimplementasikan konsep 
parameterized testing menggunakan JUnit 5. Dengan berbagai sumber data seperti
@ValueSource, file CSV (@CsvFileSource), dan @MethodSource, pengujian dapat 
dijalankan berkali‑kali dengan berbagai input tanpa mengulang kode. 

Beberapa hal penting yang diperoleh dari praktikum:
\begin{itemize}
    \item Penggunaan anotasi @ParameterizedTest memudahkan pengujian banyak kasus.
    \item Sumber data fleksibel (literal, CSV, objek) mendukung skenario pengujian
          yang kompleks.
    \item Penanganan exception seperti InsufficientFundsException dan
          validasi IllegalArgumentException dapat diuji secara terstruktur.
    \item Refaktor kelas pendukung (Owner, Card, Wallet) meningkatkan keterbacaan 
          dan kesiapan untuk menerima parameter terperinci.
\end{itemize}

Secara keseluruhan, praktikum ini berhasil memenuhi tujuan pembelajaran dengan
menunjukkan bahwa metode pengujian terparameterisasi efektif untuk meningkatkan
cakupan dan meminimalkan duplikasi. Mahasiswa diharapkan dapat menerapkan 
pendekatan serupa pada proyek pengujian berikutnya.


\section{Daftar Pustaka}
\begin{enumerate}[label={[\arabic*]}]
    \item Medium - All You Need to Know About Parameterized Tests \url{https://medium.com/@torresgmarina/all-you-need-to-know-about-parameterized-tests-493af4705508}
    \item JUnit - Parameterized Classes and Tests \url{https://docs.junit.org/6.0.3/writing-tests/parameterized-classes-and-tests.html}
    \item Baeldung - Parameterized Tests in JUnit 5 \url{https://www.baeldung.com/parameterized-tests-junit-5 }
\end{enumerate}

\end{document}